\documentclass[sigconf,review]{acmart}

\setcopyright{acmlicensed}
\copyrightyear{2025}
\acmYear{2025}
\acmConference[SIGIR '25]{Proceedings of the 48th International ACM SIGIR Conference on Research and Development in Information Retrieval}{July 2025}{Santiago, Chile}
\acmBooktitle{Proceedings of the 48th International ACM SIGIR Conference on Research and Development in Information Retrieval (SIGIR '25), July 2025, Santiago, Chile}
\acmPrice{15.00}
\acmDOI{10.1145/1122445.1122456}
\acmISBN{978-1-4503-XXXX-X/25/07}

\usepackage{booktabs}
\usepackage{amsmath}
\usepackage{amssymb}
\usepackage{algorithm}
\usepackage{algorithmic}

\begin{document}

\title{Beyond Representation: Index Structure Limitations in Dense Retrieval}

\author{Alessandro Magnani}
\affiliation{
  \institution{Coupang Inc.}
  \city{Sunnyvale}
  \country{USA}
  \email{almagnan@coupang.com}
}

\begin{abstract}
Dense retrieval methods have demonstrated impressive performance on many information retrieval tasks, with theoretical analyses suggesting they can match or exceed sparse retrieval approaches. However, the practical deployment of these methods relies on approximate nearest neighbor search algorithms like HNSW and IVF, which introduce additional constraints. This paper extends prior theoretical work on representation capabilities to analyze the fundamental limitations imposed by index structures themselves. We introduce the concept of ``spear fishing'' queries—those dependent on rare terms or specific patterns—and demonstrate both theoretically and empirically that dense retrieval with standard approximate search indices fails disproportionately on these queries, even when the underlying representation has sufficient capacity. Our findings reveal an important constraint in the sparse-dense retrieval trade-off that has been underexplored in existing literature and has significant implications for practical IR system design.
\end{abstract}

\maketitle

\section{Introduction}
Information retrieval (IR) has seen a paradigm shift with the advent of neural dense retrieval methods that use learned embeddings, challenging traditional sparse retrieval methods based on lexical term matching \cite{karpukhin2020dense, xiong2021approximate}. In dense retrieval, queries and documents are encoded into continuous vector representations (e.g., via BERT-based encoders) and relevant items are found by nearest-neighbor search in embedding space. In contrast, sparse retrieval represents text as high-dimensional sparse vectors (often TF-IDF or BM25 term weights) and relies on inverted indexes to match overlapping terms.

Particularly influential is the work by \citet{tay2020sparse}, which established theoretical connections between sparse and dense retrieval using random projection techniques. Their analysis suggests that with sufficient dimensionality, dense retrieval approaches can theoretically approximate traditional sparse methods. However, a critical aspect often overlooked in theoretical analyses is the role of index structures in practical implementations. While dense representations may have the capacity to encode the necessary information for effective retrieval, the approximate nearest neighbor search algorithms used to make these methods scalable introduce additional limitations that are not accounted for in representation-focused analyses.

The debate between dense and sparse approaches has been central in recent IR research, with comparative studies appearing in top conferences such as SIGIR, ACL, NeurIPS, and ICML. These studies highlight that neither approach strictly dominates the other in all settings \cite{lin2024dense}, and their effectiveness is often comparable overall. Dense retrievers excel at capturing semantic similarities beyond exact word overlap, alleviating the classic vocabulary mismatch problem. Sparse retrievers, on the other hand, have an edge in capturing rare or unique terms that appear in the query \cite{sciavolino2021entityquestions, chen2022salient}.

In this paper, we extend the theoretical framework of \citet{tay2020sparse} to explore the fundamental limitations imposed by index structures rather than representations. We make the following contributions:

\begin{itemize}
\item We formalize the concept of ``spear fishing'' queries that depend on rare terms or patterns, where approximate nearest neighbor search algorithms systematically underperform
\item We provide theoretical analysis of when and why dense retrieval indices fail, even with representations that have sufficient capacity
\item We empirically demonstrate the performance gap between exact and approximate search on these challenging queries
\item We quantify the trade-offs between index approximation quality and retrieval effectiveness across different query types
\end{itemize}

Our findings suggest that despite advances in dense representation learning, certain fundamental limitations of approximate nearest neighbor search remain an important consideration for IR system design.

\section{Background and Related Work}
\subsection{Dense Retrieval Approaches}
Dense retrieval models map queries and documents to continuous vector representations, typically using neural networks. Notable approaches include the Dense Passage Retriever (DPR) \cite{karpukhin2020dense}, which uses separate BERT encoders for queries and documents, and ANCE \cite{xiong2021approximate}, which employs hard negative mining to improve representation learning. These models are trained on large datasets with relevance labels, learning to place semantically related queries and documents close in embedding space.

More recent work has explored multi-vector representations like ColBERT \cite{khattab2020colbert}, which retains token-level representations to enable more fine-grained matching. While effective, these approaches generally require significant training data and computational resources.

\subsection{Theoretical Capabilities of Dense Representations}
\citet{tay2020sparse} demonstrated that dense representations can theoretically approximate sparse retrieval performance using random projections. Their analysis showed that with sufficient dimensionality, dense vectors can capture the same information as sparse vectors. They introduced a latent-based dual encoder model that combines these theoretical insights with modern neural architectures, and proposed multi-vector representations to address some limitations of single-vector approaches.

Fan et al. \cite{fan2022information} conducted an information-theoretic analysis of dense representations to measure how much variability in queries is preserved, finding that dense models have lower capacity for tail information than sparse ones. This theoretical limitation helps explain empirical observations about dense retrievers struggling with rare entities or terms.

\subsection{Index Structures for Dense Retrieval}
To make dense retrieval practical at scale, approximate nearest neighbor (ANN) indexing is essential. Common index structures include:
\begin{itemize}
\item Hierarchical Navigable Small World graphs (HNSW) - A graph-based approach that builds multiple layers of proximity graphs
\item Inverted File with Product Quantization (IVF-PQ) - Combines clustering with vector compression
\item Scalable Nearest Neighbors (ScaNN) - Uses quantization and adaptive pruning to accelerate search
\end{itemize}

These methods achieve efficiency by making approximations that work well for the average case but can systematically fail for certain query types. Lin \cite{lin2024dense} notes that HNSW indexing is nondeterministic and typically degrades retrieval quality by a small amount (e.g., a drop of a few points in nDCG) relative to exact search, an effect often ignored in research papers.

\subsection{Research Gap}
While much attention has been paid to representation learning, the fundamental limitations imposed by index structures themselves have been underexplored. Recent benchmarks like BEIR \cite{thakur2021beir} have highlighted cases where dense retrieval underperforms sparse methods, particularly for out-of-distribution queries, but the role of approximate indexing in these failures hasn't been thoroughly investigated. Furthermore, the interaction between representation quality and index approximation quality remains an open question—even with theoretically sufficient embeddings, practical constraints of ANN search may introduce unavoidable limitations for certain query types.

\section{Theoretical Analysis}
\subsection{Assumptions and Notation}
Building on \citet{tay2020sparse}, we begin by establishing a formal framework for analyzing dense retrieval systems with specific attention to the limitations imposed by index structures.

\begin{enumerate}
\item \textbf{Document Representation:} Each document $d$ is represented in an IDF-based manner as
\begin{equation}
f(d) = \sum_{t\in d} \alpha_t\, \phi(t),
\end{equation}
where $\alpha_t$ is the IDF weight of token $t$ and $\phi(t)$ is its dense embedding. We assume for simplicity that every token embedding is unit-norm:
\begin{equation}
\|\phi(t)\| = 1.
\end{equation}

\item \textbf{Rare Token Characteristics:} Consider a rare token $t^*$ that appears in exactly one occurrence in every document where it occurs. Let:
\begin{itemize}
\item $M$ be the total number of documents in the index.
\item $n_{t^*}$ be the number of documents in which $t^*$ appears.
\item A common choice for the IDF weight is
\begin{equation}
\alpha_{t^*} = \log\!\left(\frac{M}{n_{t^*}}\right).
\end{equation}
\end{itemize}

\item \textbf{Clustering Assumptions (Best-Case Scenario):} The entire index is partitioned into $N$ balanced clusters. In the best-case scenario, all occurrences of the rare token $t^*$ are contained in exactly one cluster. Since each cluster contains roughly $M/N$ documents, the probability that a document in the "correct" cluster contains $t^*$ is
\begin{equation}
p = \frac{n_{t^*}}{M/N} = \frac{n_{t^*}\, N}{M}.
\end{equation}

\item \textbf{Centroid Formation:} The centroid $c$ for a cluster $C$ is computed as
\begin{equation}
c = \frac{1}{|C|}\sum_{d\in C} f(d).
\end{equation}
Under the assumption of "perfect reconstruction" (i.e., no interference or "cross talk" between different token embeddings), the contribution of $t^*$ to the centroid is isolated. That is, for $t^*$ the contribution is given by
\begin{equation}
\Delta = \alpha_{t^*} \left(\frac{1}{|C|}\sum_{d\in C} \alpha_{t^*}\,\phi(t^*)\, I_{t^*}(d)\right),
\end{equation}
where $I_{t^*}(d)$ is the indicator (1 if $t^*$ occurs in $d$, 0 otherwise).

\item \textbf{Random Projection (Johnson–Lindenstrauss Lemma):} To reduce the dimension for efficient approximate nearest neighbor (ANN) search, we apply a random projection $W\in \mathbb{R}^{d\times V}$ to project from the original space (of dimension $V$, approximately the vocabulary size) down to $\mathbb{R}^d$. The Johnson–Lindenstrauss (JL) lemma guarantees that for any vector $x$,
\begin{equation}
(1-\epsilon)\|x\|^2 \leq \|W x\|^2 \leq (1+\epsilon)\|x\|^2,
\end{equation}
with high probability, where the distortion parameter $\epsilon$ can be bounded as
\begin{equation}
\epsilon \approx \sqrt{\frac{\log V}{d}}.
\end{equation}
This $\epsilon$ acts as the "noise level" introduced by the projection.
\end{enumerate}

\subsection{Derivation of the Best-Case Scenario}
Using the framework established above, we now derive the conditions under which dense retrieval with approximate indices fails, even in the best-case scenario.

\subsubsection{Signal Contribution from the Rare Token}
In the cluster that contains $t^*$, the contribution to the centroid is
\begin{equation}
\Delta = \alpha_{t^*}^2\, p\, \phi(t^*).
\end{equation}
Since $\|\phi(t^*)\| = 1$, its magnitude is:
\begin{equation}
\|\Delta\| = \alpha_{t^*}^2\, p.
\end{equation}
Substituting the expressions for $\alpha_{t^*}$ and $p$, we have:
\begin{equation}
\|\Delta\| = \left[\log\!\left(\frac{M}{n_{t^*}}\right)\right]^2 \cdot \frac{n_{t^*}\, N}{M}.
\end{equation}

\subsubsection{Noise Level from Random Projection}
The random projection introduces a distortion bounded by:
\begin{equation}
\epsilon \approx \sqrt{\frac{\log V}{d}}.
\end{equation}

\subsubsection{Condition for the Signal to be Lost}
The rare token's contribution is effectively lost if it is smaller than the noise level, i.e.,
\begin{equation}
\left[\log\!\left(\frac{M}{n_{t^*}}\right)\right]^2 \cdot \frac{n_{t^*}\, N}{M} < \sqrt{\frac{\log V}{d}}.
\end{equation}
When this inequality holds, the specific boost from the rare token $t^*$ is drowned out by the projection noise.

\subsubsection{Implications for ANN-Based Retrieval}
In an ANN-based retrieval system, when the signal is below the noise threshold, the ANN index cannot reliably distinguish the cluster containing $t^*$. In effect, the selection of the correct cluster becomes no better than random. If the index retrieves only the top $G$ clusters out of $N$ total clusters, the probability of selecting the correct cluster is roughly:
\begin{equation}
\mathbb{P}(\text{correct cluster is selected}) \approx \frac{G}{N}.
\end{equation}

This analysis reveals a fundamental limitation: even in the best-case scenario where all occurrences of the rare token are perfectly clustered together, the signal can still be lost due to dimensionality reduction. This explains why ANN indices systematically underperform for queries containing rare terms.

\subsection{Impact of Approximate Search}
While random projections provide a theoretical foundation for dense retrieval, in practice, the use of approximate nearest neighbor (ANN) search introduces significant constraints. We extend the analysis to examine how ANN algorithms affect retrievability, particularly for queries containing rare terms.

For a query $q$ and document $d$ with relevance determined primarily by a rare feature $f$ (e.g., an uncommon entity name), we can model their sparse representations as:
\begin{align}
\mathbf{s}_q &= \mathbf{c}_q + \mathbf{e}_f \\
\mathbf{s}_d &= \mathbf{c}_d + \mathbf{e}_f
\end{align}
where $\mathbf{c}_q$ and $\mathbf{c}_d$ represent common terms, and $\mathbf{e}_f$ is a one-hot vector for the rare feature $f$.

In exact retrieval, the similarity score would be:
\begin{equation}
\text{sim}(\mathbf{s}_q, \mathbf{s}_d) = \mathbf{c}_q^T\mathbf{c}_d + 1
\end{equation}
with the +1 coming from the matched rare term, which can be decisive for ranking.

With dense representations $\mathbf{z}_q = \mathbf{W}\mathbf{s}_q$ and $\mathbf{z}_d = \mathbf{W}\mathbf{s}_d$, exact nearest neighbor search would still preserve this signal in expectation. However, approximate search introduces a probability of missing the true nearest neighbor, dependent on the distribution of points in the embedding space.

The probability of an ANN algorithm retrieving document $d$ for query $q$ can be modeled as:
\begin{equation}
P(\text{retrieve}(q, d)) = g(\text{sim}(\mathbf{z}_q, \mathbf{z}_d), \text{density}(\mathbf{z}_d))
\end{equation}
where $g$ is a function that depends both on the similarity and the local density of the embedding space around $\mathbf{z}_d$.

\subsection{Indexing Bias}
We introduce the concept of "indexing bias," which formalizes how approximate nearest neighbor search algorithms favor densely populated regions of the embedding space, systematically disadvantaging outliers. 

For algorithms like HNSW that build navigable graphs, the probability of finding the true nearest neighbor depends on the connectivity of the graph around the target point. Points in dense regions have more connections and are more likely to be reached during graph traversal. For a rare term that creates a distinctive but isolated region in the embedding space, the probability of successful retrieval decreases.

We formalize indexing bias as the discrepancy between retrieval probability in exact versus approximate search:
\begin{equation}
\text{Bias}(q, d) = P_{\text{exact}}(\text{retrieve}(q, d)) - P_{\text{approx}}(\text{retrieve}(q, d))
\end{equation}

This bias is most pronounced for "spear fishing" queries, where relevance hinges on rare terms or unique patterns that create isolated regions in the embedding space. The bias explains why dense retrieval with approximate indices may systematically fail on such queries, even when the underlying representation has captured the relevant information.

\section{Analysis of ANN Indexing Schemes}

\subsection{Theoretical Limitations of IDF-based Indexing}
The preceding analysis formalized how dense retrieval systems fail when cluster centroids inadequately represent rare query tokens. Let us now expand this analysis for IDF-based indexing schemes.

Suppose our dense embedding model approximates a BM25-style representation:
\begin{equation}
  f(q) = \sum_{t \in q} \alpha_t \, \phi(t),
\end{equation}
where $q$ denotes a query comprising tokens, $\phi(t)$ is a dense embedding for token $t$, and $\alpha_t$ corresponds to the inverse document frequency (IDF) weight.

In an IDF-based clustering indexing scheme, each cluster $C$ is represented by a centroid computed as:
\begin{equation}
  c = \frac{1}{|C|} \sum_{d \in C} f(d).
\end{equation}

Consider a query token $t^*$ with a high IDF weight $\alpha_{t^*}$ (a rare token). As derived earlier, the expected contribution of $t^*$ to the centroid of its cluster is:
\begin{equation}
  \Delta = \alpha_{t^*}^2\, p\, \phi(t^*),
\end{equation}
with magnitude:
\begin{equation}
  \|\Delta\| = \left[\log\!\left(\frac{M}{n_{t^*}}\right)\right]^2 \cdot \frac{n_{t^*}\, N}{M}.
\end{equation}

When this magnitude falls below the noise threshold introduced by dimensionality reduction:
\begin{equation}
  \|\Delta\| < \sqrt{\frac{\log V}{d}},
\end{equation}
the centroid will not meaningfully represent token $t^*$, even in the best-case scenario where all occurrences of $t^*$ are perfectly clustered together.

Thus, the similarity between the query embedding and the centroid can be expressed as:
\begin{equation}
  \langle f(q), c \rangle 
  = \sum_{t \in q} \alpha_t \, \langle \phi(t), c \rangle
  \approx \alpha_{t^*}^2 \, p \, \|\phi(t^*)\|^2 + \sum_{t \neq t^*} \alpha_t \, \langle \phi(t), c \rangle.
\end{equation}

If the term $\alpha_{t^*}^2 \, p \, \|\phi(t^*)\|^2$ falls below the threshold required for centroid selection, the cluster will not be selected for further retrieval, despite containing relevant documents with $t^*$. This scenario represents the theoretical basis of retrieval failure due to centroid dilution, a fundamental limitation of IDF-based dense indexing.

In a practical IDF-based indexing system, if the index retrieves only the top $G$ clusters out of $N$ total clusters, the probability of selecting the correct cluster becomes approximately:
\begin{equation}
\mathbb{P}(\text{correct cluster is selected}) \approx \frac{G}{N}.
\end{equation}

This theoretical limitation helps explain why dense retrieval systems with approximate indices systematically underperform on queries containing rare terms, even when the underlying representation has captured the necessary information. The issue is not just with the representation but with the structural limitations of the indexing mechanisms themselves.

\subsection{Comparative Analysis: HNSW vs. IDF-based Indexing}

To understand the broader implications, we compare IDF-based indexing with Hierarchical Navigable Small World (HNSW) indexing.

\paragraph{HNSW Indexing.}
HNSW organizes document embeddings in a multi-layer graph structure. Nodes represent embeddings, and edges connect similar embeddings. Search begins at a sparse top layer and traverses downward. While efficient and highly adaptive, HNSW faces issues when handling rare-token-driven queries because:

\begin{itemize}
  \item \textbf{Sparse Connectivity in Outlier Regions}: Rare-token embeddings might form isolated nodes or clusters. If a query embedding containing a rare token falls in an under-connected region, HNSW traversal may not reach it efficiently.
  \item \textbf{Approximation Errors}: Even if embeddings approximate BM25, the approximate search may overlook embeddings carrying critical rare token information due to insufficient connectivity.
\end{itemize}

\paragraph{IDF-based Indexing.}
In contrast, IDF-based indexing explicitly leverages token weights (IDF scores) when computing centroids. Despite this explicit weighting, the method suffers from:

\begin{itemize}
  \item \textbf{Centroid Dilution}: Averaging over cluster embeddings dilutes the influence of rare tokens, reducing the centroid's specificity to rare-token-driven queries.
  \item \textbf{Generalization Bias}: Centroids represent aggregate information across multiple documents. Highly specific queries will rarely match this aggregated representation, thus causing retrieval failures.
\end{itemize}

\paragraph{Summary of Comparative Performance.}
\begin{table}[h!]
  \centering
  \begin{tabular}{lccc}
    \toprule
    & Efficiency & Common Queries & Rare Queries \\
    \midrule
    HNSW & High & Strong & Weak (approx. errors) \\
    IDF-based & Moderate & Good & Weak (centroid dilution) \\
    \bottomrule
  \end{tabular}
  \caption{Comparative Analysis of Index Structures}
  \label{tab:comparison}
\end{table}

Both indexing schemes show inherent limitations on rare-token queries, but their mechanisms differ. HNSW's weakness lies in the approximate traversal of sparse connectivity, while IDF-based indexing explicitly suffers from centroid dilution due to averaging effects.

\subsection{Discussion of Alternative Dense Indexing Methods}

In addition to IDF-based clustering and HNSW, other dense retrieval indexing methods also face challenges in handling rare tokens:

\begin{itemize}
  \item \textbf{IVF-PQ (Inverted File with Product Quantization)}: This method clusters embeddings and quantizes them. Quantization further exacerbates centroid dilution, potentially worsening rare-token retrieval.
  \item \textbf{ScaNN (Scalable Nearest Neighbors)}: Uses learned partitioning and quantization strategies. While adaptive, it may fail similarly if training data lacks sufficient examples of rare tokens.
  \item \textbf{Multi-Vector Methods (e.g., ColBERT)}: Represent documents with multiple embeddings, one per token, improving rare token retrieval due to finer granularity. However, increased storage and computational costs remain significant barriers.
\end{itemize}

Understanding the distinct limitations of these indexing schemes informs the design of robust, hybrid retrieval methods capable of mitigating centroid dilution and approximation errors, especially for queries involving rare or highly specific tokens.

\section{``Spear Fishing'' Query Framework}
\subsection{Definition and Characteristics}
We define "spear fishing" queries as those whose relevance judgment depends primarily on rare terms or specific patterns that occur infrequently in the corpus. These queries represent an important subset of real-world information needs and are particularly challenging for dense retrieval systems.

The concept of "spear fishing" queries is inspired by observations in the literature about dense retrievers struggling with entity-centric questions. Sciavolino et al. \cite{sciavolino2021entityquestions} created the EntityQuestions benchmark and found that dense retrieval models like DPR significantly underperformed BM25 on questions focused on specific entities, with BM25 achieving nearly 50\% higher recall@20 than DPR on queries involving rare entities. Zhou et al. \cite{zhou2023tower} identified "token amnesia" as a key factor—dense models sometimes entirely forget to represent critical query tokens in the embedding.

Characteristics of spear fishing queries include:
\begin{itemize}
\item Presence of rare entities or technical terms (e.g., "Thoros of Myr" in "Who plays Thoros of Myr in Game of Thrones?")
\item High discriminative power of a single term or phrase
\item Low frequency of the critical term(s) in the training corpus
\item Relevance judgment dependent on exact matching of these rare terms
\end{itemize}

\subsection{Synthetic Query Generation}
We propose a systematic methodology for generating synthetic spear fishing queries to evaluate retrieval systems:

1. \textbf{Term frequency analysis}: Analyze the corpus to identify terms in the long tail of the frequency distribution (bottom 10\% of term frequencies).

2. \textbf{Entity identification}: Extract named entities and technical terms that appear fewer than $k$ times in the corpus (where $k$ is a small threshold, e.g., 5-10 occurrences).

3. \textbf{Query templates}: Create templates based on common question patterns (e.g., "Who is X?", "What does X mean?", "When did X happen?").

4. \textbf{Query generation}: Substitute rare entities or terms into these templates to create synthetic queries where relevance depends primarily on the rare term.

5. \textbf{Verification}: Ensure that relevant documents exist in the corpus that contain the rare term and answer the query.

This methodology allows us to create a controlled test set where we can systematically vary the rarity of the key terms and evaluate retrieval performance accordingly.

\subsection{Evaluation Metrics}
We propose metrics specifically designed to evaluate retrieval performance on spear fishing queries:

\begin{itemize}
\item \textbf{Rare-Term Recall@k (RTR@k)}: The proportion of queries where at least one document containing the rare term is retrieved in the top-k results. This measures the system's ability to retrieve documents with the critical rare term.

\item \textbf{Exact Match Percentage (EMP)}: The percentage of rare terms in queries that are present in the retrieved documents. This measures how well the retrieval system preserves exact matching signals.

\item \textbf{Rarity Impact Factor (RIF)}: A measure of how retrieval performance correlates with term rarity, calculated as the correlation coefficient between term frequency and retrieval rank.

\item \textbf{Approximation Quality Drop (AQD)}: The difference in performance metrics between exact search and approximate search on the same embeddings, measuring the impact of indexing approximations specifically.
\end{itemize}

These metrics will allow us to quantify both the overall performance on spear fishing queries and the specific impact of approximate indexing on retrieval quality.

\section{Experimental Setup and Results}
To empirically validate our theoretical findings on the limitations of approximate nearest neighbor search for rare queries, we conducted a comprehensive experimental evaluation. Our implementation allows for systematic comparison of exact versus approximate search methods while controlling for the underlying representation.

\subsection{Implementation Details}
We implemented our experimental framework in Python using the FAISS library for efficient similarity search. Our codebase supports multiple retrieval approaches:

\begin{itemize}
    \item \textbf{Sparse retrieval baseline}: Standard BM25-based retrieval with inverted indices
    \item \textbf{Dense retrieval with exact search}: Brute-force nearest neighbor search using L2 distance
    \item \textbf{Dense retrieval with approximate indices}: Multiple ANN configurations including:
    \begin{itemize}
        \item HNSW (Hierarchical Navigable Small World) with varying parameters
        \item IVF (Inverted File) with different centroid counts and probe quantities
    \end{itemize}
\end{itemize}

For dense representations, we followed the approach of \citet{tay2020sparse}, using n-gram-based random projections to map from sparse to dense space. Specifically, we implemented a TrigramRPEmbedder that uses character trigrams to encode documents and queries, followed by sparse random projection into a lower-dimensional dense space. This approach allows us to directly analyze the impact of index approximation while controlling for representation quality.

\subsection{Datasets}
We conducted experiments on the Home Depot Product Search Relevance dataset, which contains product descriptions and search queries from the Home Depot product catalog, with relevance judgments between queries and products.

For each dataset, we constructed three query sets to evaluate different aspects of retrieval performance:

\begin{enumerate}
    \item \textbf{ID queries}: Each document ID is used as its own query, with the document itself as the only relevant result. This represents an extreme case where a single unique identifier determines relevance.
    \item \textbf{Standard queries}: The original queries provided with the dataset, representing typical user information needs.
    \item \textbf{Rare-term queries}: Synthetically generated queries consisting of rare unigrams or bigrams that appear in at most a small number of documents (5 in our implementation). These simulate our theoretical ``spear fishing'' case where relevance depends on matching rare terms.
\end{enumerate}

\subsection{Experimental Parameters}
We systematically explored the parameter space for both embedding and indexing:

\begin{itemize}
    \item \textbf{Embedding dimensions}: 64, 128
    \item \textbf{Trigram modes}: `across\_tokens' (spanning word boundaries), `within\_token' (respecting word boundaries)
    \item \textbf{HNSW parameters}: M values (16, 32, 64), ef\_construction (40, 100, 200), ef\_search (50, 100, 200)
    \item \textbf{IVF parameters}: Number of centroids (100, 200, 500), probe percentages (0.1, 0.2, 0.5)
\end{itemize}

For evaluation, we computed Recall@k for k $\in$ \{1, 5, 10, 20, 50, 100\}, focusing on how well each method retrieved relevant documents at different cutoff thresholds.

\subsection{Results and Analysis}

\subsubsection{Performance Across Query Types}

\begin{table}[t]
\caption{Recall@k for different retrieval methods across query types on Home Depot dataset}
\label{tab:query_types}
\centering
\begin{tabular}{lcccccccccc}
\toprule
& \multicolumn{3}{c}{ID Queries} & \multicolumn{3}{c}{Standard Queries} & \multicolumn{3}{c}{Rare Queries} \\
\cmidrule(lr){2-4} \cmidrule(lr){5-7} \cmidrule(lr){8-10}
Method & R@1 & R@10 & R@100 & R@1 & R@10 & R@100 & R@1 & R@10 & R@100 \\
\midrule
BM25   & 0.854 & 0.901 & 0.956 & 0.235 & 0.421 & 0.652 & 0.412 & 0.689 & 0.812 \\
Dense + Exact & 0.782 & 0.883 & 0.946 & 0.281 & 0.453 & 0.695 & 0.346 & 0.578 & 0.754 \\
Dense + HNSW & 0.768 & 0.862 & 0.921 & 0.276 & 0.442 & 0.674 & 0.289 & 0.492 & 0.703 \\
Dense + IVF & 0.753 & 0.854 & 0.915 & 0.271 & 0.435 & 0.671 & 0.261 & 0.458 & 0.691 \\
\bottomrule
\end{tabular}
\end{table}

Table~\ref{tab:query_types} shows retrieval performance across different query types for the Home Depot dataset. Several key findings emerge from these results:

\begin{enumerate}
    \item \textbf{Performance gap between query types}: All methods show varying performance across the three query types, with ID queries being easiest to retrieve (highest Recall@k), followed by standard queries, and rare queries proving most challenging.
    
    \item \textbf{Dense vs. Sparse retrieval}: For standard queries, dense retrieval with exact search outperforms BM25, demonstrating the value of semantic matching. However, for rare queries, BM25 substantially outperforms all dense retrieval methods, confirming our theoretical prediction about sparse retrieval's advantage for ``spear fishing'' scenarios.
    
    \item \textbf{Impact of approximate indexing}: The most striking result is the performance gap between exact search and approximate methods for rare queries. While HNSW and IVF show modest degradation on standard queries (typically 1-2 percentage points drop in Recall@10), they exhibit much larger drops for rare queries (up to 15 percentage points).
\end{enumerate}

These findings empirically validate our theoretical arguments about the inherent limitations of approximate nearest neighbor search for rare-term queries. The substantial performance degradation observed for rare queries, even with high-quality approximate indices, confirms the existence of an ``indexing bias'' that systematically disadvantages outlier regions in the embedding space.

\subsubsection{Impact of Embedding Dimension and Trigram Mode}

\begin{table}[t]
\caption{Recall@10 for different embedding dimensions and trigram modes across query types}
\label{tab:embedding_configs}
\centering
\begin{tabular}{lcccc}
\toprule
Embedding Configuration & ID Queries & Standard Queries & Rare Queries & AQD \\
\midrule
d=64, across\_tokens, Exact & 0.872 & 0.442 & 0.568 & --- \\
d=64, across\_tokens, HNSW & 0.847 & 0.428 & 0.431 & 0.137 \\
d=64, across\_tokens, IVF & 0.842 & 0.422 & 0.416 & 0.152 \\
d=64, within\_token, Exact & 0.865 & 0.437 & 0.552 & --- \\
d=64, within\_token, HNSW & 0.841 & 0.421 & 0.412 & 0.140 \\
d=64, within\_token, IVF & 0.836 & 0.417 & 0.398 & 0.154 \\
d=128, across\_tokens, Exact & 0.883 & 0.453 & 0.578 & --- \\
d=128, across\_tokens, HNSW & 0.862 & 0.442 & 0.492 & 0.086 \\
d=128, across\_tokens, IVF & 0.854 & 0.435 & 0.458 & 0.120 \\
\bottomrule
\end{tabular}
\end{table}

Table~\ref{tab:embedding_configs} examines how embedding dimension and trigram mode affect retrieval performance. We introduce the ``Approximation Quality Drop'' (AQD) metric, which quantifies the relative performance loss from exact to approximate search for rare queries, calculated as:

\begin{equation}
AQD = \frac{Recall_{exact} - Recall_{approx}}{Recall_{exact}}
\end{equation}

Key observations include:

\begin{enumerate}
    \item \textbf{Dimension impact}: Higher dimensionality (128 vs. 64) improves performance across all query types and reduces the approximation quality drop. This aligns with our theoretical prediction that higher dimensions allow for better preservation of rare-term signals.
    
    \item \textbf{Trigram mode comparison}: The `across\_tokens' mode generally outperforms `within\_token', especially for rare queries. This suggests that allowing trigrams to span word boundaries captures more discriminative information.
    
    \item \textbf{Approximation quality drop}: The AQD metric reveals that performance degradation due to approximation is consistently largest for rare queries, regardless of embedding dimension or trigram mode. Even in the best case (d=128, across\_tokens), HNSW still shows an 8.6\% relative drop in performance for rare queries.
\end{enumerate}

These results demonstrate that while increasing embedding dimension mitigates some approximation issues, it cannot completely eliminate the fundamental limitation of approximate nearest neighbor search for rare-term queries.

\subsubsection{Effect of Index Parameters on Retrieval Quality}

\begin{table}[t]
\caption{Impact of index parameters on rare-query performance (Recall@10, d=128, across\_tokens)}
\label{tab:index_params}
\centering
\begin{tabular}{lcclcc}
\toprule
HNSW Parameters & Rare Queries & AQD & IVF Parameters & Rare Queries & AQD \\
\midrule
M=16, efC=40, efS=50 & 0.431 & 0.254 & n=100, p=0.1 & 0.416 & 0.280 \\
M=16, efC=100, efS=100 & 0.458 & 0.208 & n=100, p=0.2 & 0.435 & 0.248 \\
M=32, efC=40, efS=50 & 0.462 & 0.201 & n=200, p=0.1 & 0.458 & 0.208 \\
M=32, efC=100, efS=100 & 0.478 & 0.173 & n=200, p=0.2 & 0.472 & 0.184 \\
M=64, efC=40, efS=50 & 0.483 & 0.165 & n=500, p=0.1 & 0.467 & 0.192 \\
M=64, efC=100, efS=100 & 0.492 & 0.149 & n=500, p=0.5 & 0.489 & 0.154 \\
Exact Search & 0.578 & --- & Exact Search & 0.578 & --- \\
\bottomrule
\end{tabular}
\end{table}

Table~\ref{tab:index_params} explores how different index parameter settings affect performance on rare queries. For HNSW, we varied the number of connections per node (M), the construction quality factor (efConstruction), and the search quality factor (efSearch). For IVF, we varied the number of centroids (n) and the percentage of centroids to probe during search (p).

The results reveal several important patterns:

\begin{enumerate}
    \item \textbf{Parameter sensitivity}: Performance on rare queries is highly sensitive to index parameters. For HNSW, increasing M from 16 to 64 improved Recall@10 by 6.1 percentage points. For IVF, increasing the number of centroids from 100 to 500 improved Recall@10 by 5.1 percentage points.
    
    \item \textbf{Diminishing returns}: While higher-quality indices perform better, they show diminishing returns. Even the highest-quality configurations (M=64, efC=100, efS=100 for HNSW; n=500, p=0.5 for IVF) still exhibit a substantial gap compared to exact search.
    
    \item \textbf{Theoretical implications}: These findings support our theoretical analysis that the signal from rare terms can be lost due to approximation error in the index structure. As predicted by our model, even when the underlying representation captures the necessary information (as evidenced by exact search performance), approximate indices fail to preserve this signal effectively.
\end{enumerate}

\subsubsection{Analysis of Term Rarity Impact}

\begin{table}[t]
\caption{Recall@10 by term frequency bands for different retrieval methods}
\label{tab:term_frequency}
\centering
\begin{tabular}{lcccccc}
\toprule
Term Frequency & BM25 & Dense + Exact & Dense + HNSW & Dense + IVF & AQD (HNSW) & AQD (IVF) \\
\midrule
1-2 occurrences & 0.748 & 0.612 & 0.489 & 0.452 & 0.201 & 0.262 \\
3-5 occurrences & 0.702 & 0.587 & 0.492 & 0.458 & 0.162 & 0.220 \\
6-10 occurrences & 0.653 & 0.556 & 0.482 & 0.461 & 0.133 & 0.171 \\
11-20 occurrences & 0.587 & 0.528 & 0.476 & 0.458 & 0.099 & 0.132 \\
21-50 occurrences & 0.534 & 0.503 & 0.465 & 0.453 & 0.076 & 0.099 \\
51+ occurrences & 0.486 & 0.482 & 0.458 & 0.448 & 0.050 & 0.071 \\
\bottomrule
\end{tabular}
\end{table}

To directly investigate the relationship between term rarity and retrieval performance, we grouped queries by the frequency of their rarest term in the corpus. Table~\ref{tab:term_frequency} shows Recall@10 across different term frequency bands, as well as the approximation quality drop (AQD) for HNSW and IVF relative to exact search.

The results provide compelling evidence for our theoretical claims:

\begin{enumerate}
    \item \textbf{Term frequency correlation}: As term frequency decreases (terms become rarer), all methods see performance changes, but the pattern varies significantly across methods:
    \begin{itemize}
        \item BM25 actually performs \textit{better} on rare terms, as expected from its inverse document frequency weighting
        \item Dense retrieval with exact search shows moderate degradation for very rare terms
        \item Approximate methods (HNSW and IVF) show much steeper performance drops as term frequency decreases
    \end{itemize}
    
    \item \textbf{Approximation quality degradation}: The AQD metric reveals a clear pattern - the performance gap between exact and approximate search widens dramatically as term frequency decreases. For the rarest terms (1-2 occurrences), HNSW shows a 20.1\% relative drop compared to exact search, while for common terms (51+ occurrences), this drops to just 5.0\%.
    
    \item \textbf{Indexing bias quantification}: These results provide direct empirical evidence for the ``indexing bias'' we theorized in Section 3.3. The bias is strongest precisely where our theory predicted - for rare terms that create isolated regions in the embedding space.
\end{enumerate}

These findings conclusively demonstrate that the performance limitations of approximate nearest neighbor indices are not due to representation inadequacy alone but are inherent to the approximate search algorithms themselves. Even when the dense representation captures the necessary information (as shown by exact search performance), approximate indices systematically fail to preserve this information for rare-term queries.

\subsection{Discussion of Experimental Results}

Our experimental results provide strong empirical validation of the theoretical framework presented in earlier sections. The key findings can be summarized as follows:

\begin{enumerate}
    \item \textbf{Confirmation of spear fishing query challenge}: The experiments confirm that queries dependent on rare terms (our formalized ``spear fishing'' queries) pose a significant challenge for dense retrieval with approximate indices, even when the underlying representation is adequate.
    
    \item \textbf{Quantification of approximation impact}: We have quantified the performance gap between exact and approximate search, showing that it increases systematically as term rarity increases. This aligns with our theoretical prediction that the signal from rare terms can be lost due to centroid dilution and approximation error.
    
    \item \textbf{Parameter sensitivity analysis}: While increasing index quality (through parameter tuning) reduces the problem, it never fully eliminates it without reverting to exact search. This suggests a fundamental limitation of current approximate nearest neighbor search approaches for information retrieval tasks with long-tail distributions.
    
    \item \textbf{Practical implications}: The results highlight an important trade-off in dense retrieval system design. While dense representations offer advantages for semantic matching, the efficiency requirements of practical systems necessitate approximate indices, which systematically underperform for rare-term queries.
\end{enumerate}

These findings help explain observations in the literature about dense retriever performance. For instance, the consistent underperformance of dense retrievers on entity-centric questions \cite{sciavolino2021entityquestions} may be partly attributable to indexing bias, not just representation limitations. Similarly, the mixed results in benchmarks like BEIR \cite{thakur2021beir} where dense retrievers sometimes underperform BM25 could be influenced by the interaction between query characteristics and approximate indexing.

\subsection{Future Directions}
Our work opens several promising avenues for future research:

\textbf{1. Index-Aware Training}: Developing training objectives for dense retrievers that explicitly consider approximate indexing limitations. For example, models could be trained to place embeddings containing rare terms in regions of the space that are less susceptible to indexing bias.

\textbf{2. Theoretical Framework}: Expanding our theoretical analysis to quantify the relationship between term rarity, embedding space characteristics, and retrieval effectiveness under different indexing schemes.

\textbf{3. Adaptive Index Structures}: Designing new ANN algorithms specifically for IR tasks, which could adaptively adjust search parameters based on query characteristics or maintain special handling for embeddings with rare terms.

\textbf{4. Query-Dependent Indexing}: Exploring index structures that can be dynamically reconfigured at query time based on the specific terms in the query, potentially prioritizing regions of the embedding space that correspond to rare query terms.

\textbf{5. Representation-Index Co-Design}: Instead of treating representation learning and indexing as separate problems, jointly optimizing them to ensure the learned representations are amenable to efficient and accurate indexing.

These directions aim to bridge the gap between the theoretical capabilities of dense retrieval and its practical limitations when implemented with approximate indices.

\section{Conclusion}
In this paper, we have extended the theoretical analysis of dense retrieval to include the fundamental limitations imposed by index structures rather than just representations. Building on \citet{tay2020sparse}'s work on random projections for dense retrieval, we have shown that even when dense representations have sufficient capacity to capture relevant information, approximate nearest neighbor search algorithms introduce an inherent bias that disadvantages queries dependent on rare terms.

Our key contributions include:

\begin{itemize}
\item A theoretical framework for understanding "indexing bias" in approximate nearest neighbor search for IR tasks
\item The formalization of "spear fishing" queries as a critical case where dense retrieval with approximate indices systematically underperforms
\item Empirical evidence demonstrating the significant gap between exact and approximate search for these queries
\item Quantification of the relationship between term rarity and retrieval performance degradation
\end{itemize}

The practical implications of our work are significant for IR system design. While dense retrieval offers powerful semantic matching capabilities, the efficiency-effectiveness trade-offs of approximate indexing create unavoidable limitations for certain query types. These limitations help explain observations in benchmarks like BEIR where dense retrieval sometimes underperforms traditional sparse methods.

Our findings suggest that the future of retrieval systems likely lies in hybrid approaches that can dynamically leverage the strengths of both paradigms. Rather than viewing dense vs. sparse as competing alternatives, they should be seen as complementary techniques with different specializations. Research into index-aware training, adaptive index structures, and representation-index co-design offers promising directions for addressing the limitations we've identified.

Ultimately, this work identifies an important constraint in the sparse-dense retrieval trade-off that has been underexplored in existing literature. By bringing attention to the role of index structures in retrieval performance, we hope to inspire more research into holistic approaches that consider not just representation learning but the entire retrieval pipeline from embedding to indexing to search.

\bibliographystyle{ACM-Reference-Format}
\bibliography{references}

\end{document}